\section{Mecanismo de carga de los Skills}

\subsection{Arquitectura del System Prompt en Codex}

Antes de que el usuario emita su primera solicitud, Codex ya ha insertado cuatro mensajes de sistema:

\begin{enumerate}
    \item \textbf{Prompt a nivel System}: restricciones a nivel de sistema.
    \item \textbf{Mensaje Developer}: reglas de permisos y sandbox.
    \item \textbf{Mensaje User (inyectado)}: contiene \texttt{agents.md} y la descripción de todos los Skills disponibles; esta es la \textbf{primera exposición} de los Skills.
    \item \textbf{Mensaje Developer}: define el modo de colaboración (Plan mode / Agent mode).
\end{enumerate}

\subsection{Proceso de emparejamiento y carga de un Skill}

\begin{enumerate}
    \item El usuario emite una solicitud real.
    \item El Coding Agent \textbf{empareja el Skill más adecuado} a partir del query.
    \item El Agent invoca la herramienta \texttt{execute\_commands} y ejecuta el comando \texttt{cat} para cargar el contenido completo del \texttt{skill.md} correspondiente.
    \item El contenido del Skill entra al contexto, y todas las interacciones posteriores se ejecutan a partir de la descripción de ese Skill.
    \item Al terminar un turn, el usuario puede seguir iterando sobre sus necesidades.
\end{enumerate}

\begin{knowledgebox}{Los dos orígenes de agents.md}
\begin{itemize}
    \item \textbf{A nivel de sistema}: \texttt{\~{}/.codex/agents.md} (configuración global).
    \item \textbf{A nivel de proyecto}: el \texttt{agents.md} en la raíz del repositorio de código (instrucciones propias del proyecto).
\end{itemize}
Ambos se inyectan en el tercer mensaje User.
\end{knowledgebox}

\subsection{Instalación y actualización}

\begin{itemize}
    \item Claude Code instala Superpowers mediante el mecanismo de plugin.
    \item Codex deja que el Agent se instale a sí mismo a partir de lenguaje natural.
    \item Forma de actualizar: entrar al directorio de Superpowers y ejecutar \texttt{git pull}; se carga automáticamente en la siguiente sesión nueva.
    \item La versión actual admite \textbf{descubrimiento autónomo}: al escribir \texttt{/skills} se listan todos los Skills disponibles.
\end{itemize}

\subsection{Resumen de la sección}

Los Skills, mediante Prompt Injection, exponen su directorio en el tercer mensaje User de Codex; el Agent empareja dinámicamente según el query del usuario y usa el comando \texttt{cat} para cargar el \texttt{skill.md} correspondiente, inyectando así el proceso de ingeniería de software en el contexto del modelo.
